\documentclass[conference]{IEEEtran}
\usepackage[utf8]{inputenc}
\usepackage[T1]{fontenc}
\usepackage{amsmath}
\usepackage{booktabs}
\usepackage{microtype}
\usepackage[hidelinks]{hyperref}

\begin{document}
\title{The Gate, Not the Predictor?\\
Matched-Control and Endpoint-Metric Audits of\\
Cache Prefetching and LLM Speculative Decoding}
\author{\IEEEauthorblockN{Simone Jarno Casartelli and Enrico Lopedoto}
\IEEEauthorblockA{\textit{Algorithmica Solutions}\\
\texttt{simone@algorithmicasolutions.com},
\texttt{enrico@algorithmicasolutions.com}}}
\maketitle

\begin{abstract}
Speculative systems predict future work and then decide which predictions to
admit. Evaluations often change both components at once or substitute an
admission-conditioned success ratio for the machine endpoint. We present two
matched-control audits in different domains. In native ChampSim, a
dominant-stride gate raises accepted-prefetch accuracy from 10.64\% to 15.62\%
and suppresses 36.77\% of L1D prefetches, yet changes DRAM reads by only
$-0.083\%$ and loses 0.566\% geometric-mean IPC. In a separately frozen
Qwen2.5 speculative-decoding holdout, raising a draft-confidence threshold
from 0.8 to 0.95 raises accepted-token fraction from 90.47\% to 95.95\%, yet
reduces throughput by 10.37\% on all eight prompts (exact sign-test
$p=0.0078125$). The less restrictive gate is 1.0627$\times$ target-only
throughput; the higher-acceptance gate is 0.9525$\times$. Both failures have
the same cause: a ratio improves because its gate offers less work, while
absolute useful work and downstream cost follow different paths. The
contribution is a reproducible attribution and endpoint-audit method, not a
new predictor, confidence gate, or universal performance claim.
\end{abstract}

\begin{IEEEkeywords}
speculation, matched controls, cache prefetching, speculative decoding,
performance methodology, negative results
\end{IEEEkeywords}

\section{Introduction}
Speculation is a recurring systems pattern. A cheap mechanism predicts work,
a policy decides whether to admit it, and an expensive downstream machine
executes or verifies it. Cache prefetchers predict addresses; LLM draft models
predict tokens. Confidence filters suppress prefetches or end draft spans.
Both domains commonly report a success ratio conditioned on admission:
prefetch accuracy or draft-token acceptance.

That ratio is useful but not an endpoint. If a gate removes ten unsuccessful
proposals and two successful ones, its success fraction rises while the amount
of useful work falls. The downstream system may see no traffic reduction or
latency benefit. Attribution is also lost when a gated learned predictor is
compared with an ungated classical predictor: prediction and admission change
together.

We ask whether these are domain-specific mistakes or one general experimental
failure. Our contributions are:
\begin{itemize}
\item a shared propose--admit--execute decomposition and matched-control rule;
\item a latency-aware local and native ChampSim audit that falsifies a broad
neural-prefetching interpretation;
\item a resource-bounded real-model LLM holdout that isolates only the draft
confidence threshold and exhibits a predeclared acceptance/throughput
reversal; and
\item machine-verifiable artifacts that distinguish independent units from
timing repetitions and reconstruct endpoints from raw runs.
\end{itemize}

The LLM result independently reproduces a warning present in recent
speculative-decoding work. Novelty lies in the cross-domain attribution method
and paired counterexamples, not in claiming that confidence-aware drafting is
new.

\section{A common audit model}
Let $P$ be proposed work, $A\subseteq P$ admitted work, and $S\subseteq A$
successful admitted work. A common proxy is
\begin{equation}
Q=|S|/|A|.
\end{equation}
A stricter gate can raise $Q$ while reducing both $|A|$ and $|S|$. Neither
$Q$ nor $|A|$ determines endpoint $E$, because admitted work may be absorbed,
merged, arrive late, displace useful state, or cost more to generate than it
saves.

We use three rules. First, compare gates with an identical predictor and
predictors with an identical gate. Second, report absolute work at successive
machine boundaries, not only the admission-conditioned ratio. Third, end at a
user- or machine-visible metric: IPC and DRAM traffic for caches; generation
latency and tokens/s for LLM inference.

Repeated timing of one deterministic workload estimates measurement noise but
does not create new independent units. We therefore collapse repetitions
within a trace or prompt before paired inference.

\section{Case study I: cache prefetching}
The original study compared a gated online 6--32--1 MLP delta predictor with an
always-on stride predictor. It also obtained its largest speedups under
instantaneous prefetch completion. We added classical predictors with the same
admission policy, asynchronous latency, contention, finite MSHRs, an
associative cache, and demand-only shadow state. Matched controls show no MLP
advantage on random traffic and large MLP losses on several regular streams.
The apparent safety came primarily from admission, not neural prediction.

A compact dominant-delta gate remained promising locally, so we froze a native
holdout before observation: the highest-weight public DPC-3 SimPoint for each
of nine previously unseen SPEC CPU2017 programs, with 50 million warm-up and
200 million measured instructions. A runtime ChampSim module selects no L1D
prefetching, official \texttt{ip\_stride}, a gate-disabled matched port, or
gated stride from one binary. The disabled port is output-identical to official
\texttt{ip\_stride} on every recorded field.

Table~\ref{tab:cross} summarizes the result. The gate removes 36.77\% of
accepted L1D prefetches and raises accuracy by 4.98 percentage points, while
retaining 92.83\% of useful prefetches. Only 6.85\% fewer prefetch misses reach
LLC; DRAM reads barely change. Direct gated/raw geometric-mean IPC is 0.99434
(bootstrap 95\% CI 0.98445--0.999997), with four wins and five losses. The
proxy improvement does not establish bandwidth or performance benefit.

\begin{table}[t]
\centering\small
\caption{The same endpoint-metric failure in two domains.}
\label{tab:cross}
\begin{tabular}{lrr}
\toprule
metric & looser/raw & stricter/gated \\
\midrule
Prefetch accuracy & 10.64\% & 15.62\% \\
accepted L1D prefetches & 100\% & 63.23\% \\
DRAM reads & 100\% & 99.917\% \\
IPC ratio (strict/raw) & \multicolumn{2}{c}{0.99434} \\
\midrule
Draft acceptance & 90.47\% & 95.95\% \\
accepted draft tokens & 1,623 & 1,278 \\
geomean tokens/s & 27.24 & 24.41 \\
throughput ratio (strict/loose) & \multicolumn{2}{c}{0.89628} \\
\bottomrule
\end{tabular}
\end{table}

\section{Case study II: LLM speculative decoding}
\subsection{Question and controls}
Speculative decoding uses a smaller draft model to propose tokens that a target
model verifies in parallel, preserving the target distribution in the
idealized algorithm~\cite{leviathan,chen}. Static draft length can waste work
when token difficulty changes. SpecDec++, dynamic assisted generation, and
SVIP use learned or confidence/entropy signals to stop drafting
adaptively~\cite{specdecpp,hfdynamic,svip}.

This makes speculative decoding a direct test of the shared audit model. The
draft model is the predictor, confidence is the admission/stopping gate,
accepted-token fraction is the proxy, and decode throughput is the endpoint.
We hold target, draft, quantization, prompt, output count, sampling, maximum
draft length, runtime, and hardware fixed. The primary direct contrast changes
only the minimum draft-token confidence from 0.8 to 0.95.

\subsection{Frozen design}
We use official Qwen2.5-3B-Instruct as target and Qwen2.5-0.5B-Instruct as
draft, both Q4\_K\_M GGUF~\cite{qwen}. The CPU-only \texttt{llama.cpp}
runtime is pinned to commit \texttt{5854625}. Runs use four threads pinned to
four logical CPUs, one server slot, no prompt cache, greedy sampling, and
exactly 96 generated tokens. The host is a Ryzen AI 9 HX PRO 370 with 24~GiB
system memory. GPU execution is deliberately excluded.

Four HumanEval and four GSM8K prompts form untouched holdout units. Indices 10,
25, 50, and 100 in each public test set were frozen before observation. Three
timing repetitions per prompt and policy are collapsed to their median.
Configuration order is deterministically shuffled in each repeat block.

Development prompts at indices 0 and 5 exposed two facts before holdout. First,
the 0.95 threshold gave higher acceptance but lower throughput than 0.8,
forming the frozen confirmatory prediction. Second, verification batch shape
changed some greedy Q4 outputs. We therefore forced equal output-token counts,
retained output hashes as a diagnostic, and made no output-quality claim.

\subsection{Holdout result}
The 120-run matrix is complete. Raising the threshold removes 25.75\% of draft
proposals and 21.26\% of accepted draft tokens. Acceptance fraction therefore
rises by 5.48 percentage points even as absolute accepted work falls.

Throughput moves in the opposite direction. The paired geometric-mean ratio is
0.89628, and threshold 0.95 is slower on every prompt. The exact two-sided sign
test gives $p=0.0078125$ for the single predeclared contrast. Threshold 0.8 is
1.06274$\times$ target-only throughput; threshold 0.95 is 0.95250$\times$.
Static lengths two and four are also slower than target-only at 0.94381$\times$
and 0.93749$\times$. A high accepted-token fraction is thus neither necessary
nor sufficient evidence of speedup on this CPU.

\begin{table}[t]
\centering\small
\caption{Frozen LLM holdout; medians over three timings per prompt.}
\label{tab:llm}
\begin{tabular}{lrrr}
\toprule
prompt & gate .8 & gate .95 & .95/.8 \\
\midrule
GSM8K 10  & 27.06 & 25.75 & 0.9518 \\
GSM8K 25  & 26.73 & 23.65 & 0.8848 \\
GSM8K 50  & 22.80 & 21.06 & 0.9235 \\
GSM8K 100 & 27.99 & 25.56 & 0.9130 \\
HumanEval 10  & 24.47 & 21.67 & 0.8857 \\
HumanEval 25  & 25.84 & 23.40 & 0.9057 \\
HumanEval 50  & 32.39 & 26.77 & 0.8265 \\
HumanEval 100 & 32.04 & 28.34 & 0.8845 \\
\bottomrule
\end{tabular}
\end{table}

Output hashes differ across policies for six of eight prompts, yet each
policy's hash is stable across all three repetitions. Static length two matches
target-only on every prompt, while other verification shapes sometimes select
a different continuation. This is consistent with known finite-precision and
batch-shape sensitivity of greedy inference~\cite{numerics}; it does not prove
the kernel-level cause or a distributional violation. It does show that
``lossless'' should not be silently strengthened to byte-identical output for
this quantized systems configuration.

\section{Cross-domain interpretation}
The two cases are not merely metaphorical. Both gates condition the denominator
of their headline metric. In caching, the gate suppresses 36.77\% of admitted
requests and accuracy rises, but most removed work never reached DRAM or later
returned as demand traffic. In LLM decoding, the stricter gate suppresses
25.75\% of proposals and acceptance rises, but shorter verification groups
lose more batching benefit than the saved draft work repays.

The result does not imply that gates are harmful. Gate 0.8 is the only tested
LLM policy that improves aggregate throughput, and raw stride materially beats
no-prefetch in the cache study. It implies that admission policy is a separate
component whose cost must be evaluated at the endpoint.

\section{Related work}
Classical and learned prefetchers span stride, confidence, spatial, and learned
prediction~\cite{stride,spp,bingo,learning}. Our cache contribution is an
attribution audit rather than a new state-of-the-art prefetcher.

Recent LLM work already recognizes the draft-cost/acceptance trade-off.
Learning to Draft optimizes throughput rather than acceptance-length proxies
with co-adaptive policies~\cite{ltd}. Ha et al. find acceptance rate to be a
poor standalone predictor and draft cost to dominate throughput across
self-speculative methods~\cite{char}.
Our result is an independent CPU replication with a matched threshold and a
cross-domain endpoint framework, not priority for that observation.

\section{Threats to validity}
The cache holdout uses nine programs, one highest-weight SimPoint per program,
one ChampSim configuration, and no multicore contention. The LLM holdout uses
one consumer CPU, one quantized model pair, eight prompts from two datasets,
batch size one, and fixed short outputs. The threshold was chosen on four
development prompts. Neither case measures silicon energy or production
serving load.

The LLM sign test has only eight units; its small $p$ follows unanimous
direction, not broad workload coverage. Timing repetitions are not counted as
independent. Output divergence prevents a quality claim and may couple content
with later draft difficulty even under equal token counts. Future work should
add model families, Q8 or floating-point targets, GPUs, a second runtime,
natural-language workloads, long contexts, per-round timing, and energy.

\section{Reproducibility}
The cache artifact contains 36 native runs and 21,240 local rows. The LLM
artifact contains 120 raw runs. Manifests pin source, runtime, configuration,
model, prompt, binary, and trace hashes. The verifiers reconstruct every saved
aggregate and enforce the independent-unit and fixed-token invariants:
\begin{verbatim}
python -m benchmarks.verify_research
python champsim/verify_results.py \
  --results \
  results/champsim_holdout_standard
python \
  benchmarks/specdec_endpoint_audit.py \
  verify --output \
  results/specdec_holdout_qwen_cpu
\end{verbatim}

\section{Conclusion}
Across cache prefetching and LLM speculative decoding, a stricter gate improves
an admission-conditioned success ratio while worsening the endpoint. The
shared lesson is methodological: isolate prediction from admission, track
absolute work through downstream boundaries, and end the argument at IPC,
traffic, latency, or throughput. Proxy metrics explain mechanisms; they do not
prove system benefit.

\begin{thebibliography}{14}
\bibitem{stride} T.-F. Chen and J.-L. Baer, ``Effective hardware-based data
prefetching for high-performance processors,'' \emph{IEEE TC}, 1995.
\bibitem{spp} J. Kim et al., ``Path confidence based lookahead prefetching,''
\emph{MICRO}, 2016.
\bibitem{bingo} M. Bakhshalipour et al., ``Bingo spatial data prefetcher,''
\emph{HPCA}, 2019.
\bibitem{learning} M. Hashemi et al., ``Learning memory access patterns,''
\emph{ICML}, 2018.
\bibitem{leviathan} Y. Leviathan, M. Kalman, and Y. Matias, ``Fast inference
from Transformers via speculative decoding,'' \emph{ICML}, 2023.
\bibitem{chen} C. Chen et al., ``Accelerating large language model decoding
with speculative sampling,'' arXiv:2302.01318, 2023.
\bibitem{specdecpp} K. Huang, X. Guo, and M. Wang, ``SpecDec++: Boosting
speculative decoding via adaptive candidate lengths,'' \emph{COLM}, 2025.
\bibitem{hfdynamic} J. Mamou et al., ``Faster assisted generation with dynamic
speculation,'' Hugging Face, 2024.
\bibitem{svip} Z. Zhang et al., ``Draft model knows when to stop:
Self-verification speculative decoding for long-form generation,''
\emph{EMNLP}, 2025.
\bibitem{ltd} J. Zhang et al., ``Learning to draft: Adaptive speculative
decoding with reinforcement learning,'' \emph{ICLR}, 2026.
\bibitem{char} J. Ha, K. Ganesan, A. Nguyen, T. Ahmed, and A. Moshovos,
``Characterizing self-speculative decoding approaches for accelerating
LLMs,'' \emph{ICML AdaptFM Workshop}, 2026.
\bibitem{numerics} J. Yuan et al., ``Understanding and mitigating numerical
sources of nondeterminism in LLM inference,'' \emph{NeurIPS}, 2025.
\bibitem{qwen} A. Yang et al., ``Qwen2.5 technical report,''
arXiv:2412.15115, 2024.
\bibitem{champsim} N. Gober et al., ``The Championship Simulator:
Architectural simulation for education and competition,'' arXiv:2210.14324,
2022.
\end{thebibliography}
\end{document}
